Xianxia (traditional Chinese: 仙俠; simplified Chinese: 仙侠; pinyin: xiānxiá; lit. 'immortal heroes') is a genre of Chinese fantasy heavily inspired by Chinese mythology and influenced by philosophies of Taoism, Chan Buddhism, Confucianism, Chinese martial arts, classical Chinese medicine, Chinese folk religion, Chinese alchemy, and other classical elements of Chinese culture.[1] It is also closely related to the wuxia and shenmo genre, as it developed as a subgenre of wuxia and simultaneously as a subgenre of shenmo before branching out to become a distinct martial arts genre by the late 19th century.[2][3] It is often considered akin to other high fantasy genres such as the sword and sorcery stories of Western literatures and the ninja fantasy fictions in Japanese literature.
\\

Protagonists of xianxia stories are often practitioners or trainees of magic and divination — i.e. truth cultivators (修真者) or simply cultivators (修士) — seeking spiritual enlightenment, immortality and supernatural miracle powers, or else are transcendent beings called xiān (仙), who have either already ascended or have mastered such aptitudes to varying degrees. Antagonists of xianxia stories typically have similar powers, and often belong to either malevolent rival clans/schools of cultivators, or mischievous/malicious tribes or gangs of ling (fairies and sprites), yao (monsters), mo (demons), gui (ghosts and undead beings such as zombies and vampires) and similar category of inhuman sentient beings. Persons in the xianxia genre manifest superhuman talents or physics-defying superpowers such as flight/levitation, teleportation, telekinesis, divination/soul flight, shapeshifting, materializing objects and force fields, manipulation of energy and the elements, etc.
\\

Concepts from classical Chinese philosophies such as internal alchemy and external alchemy feature in this genre—deities, immortals, yaoguai, demons and ghosts all engage in meditative practices and the consumption of rare substances or creatures to improve their skills or to augment their power. Action tends to take place across multiple realms, the number of which depends on the author or the world in question, but this usually includes the immortal plane, the mortal realm, and in the underworld. The xianxia genre also tends to feature the existence of magical creatures who do not belong to either the yao or mo category, as well as supernatural artefacts capable of upending the status quo.
